\documentclass{article}
\usepackage{amsmath,amssymb,stmaryrd,amsthm,MnSymbol}
\newcommand{\pp}[2]{\frac{\partial #1}{\partial #2}} 
\newcommand{\dede}[2]{\frac{\delta #1}{\delta #2}}
\newcommand{\dd}[2]{\frac{\diff#1}{\diff#2}}
\newcommand{\dt}[1]{\diff\!#1}
\newcommand{\jump}[1]{\left[\!\!\left[ #1 \right]\!\!\right]}
\DeclareMathOperator{\Hess}{Hess}
\def\MM#1{\boldsymbol{#1}}
\DeclareMathOperator{\diff}{d}
\DeclareMathOperator{\DIV}{DIV}
\DeclareMathOperator{\D}{D}
\bibliographystyle{plain}
\newcommand{\vecx}[1]{\MM{#1}}
\newtheorem{definition}{Definition}
\newcommand{\code}[1]{{\ttfamily #1}} 
\usepackage[margin=2cm]{geometry}

\usepackage{fancybox}
\begin{document}
\title{Notes on energy conserving method for SWE}
\author{Colin}
\maketitle

\section{Poisson bracket method}

The Poisson bracket method makes use of the vector invariant form of
the velocity advection term
\[
(u\cdot \nabla)u = u\nabla\cdot u^{\perp} + \nabla\frac{|u|^2}{2},
\]
so we write the shallow water equations in the form
\begin{align}
  \pp{u}{t} + u^\perp(\nabla\cdot u^{\perp} + f) + \nabla\left(\frac{|u|^2}{2}
  - g(D+b)\right) & = 0, \\
  \pp{D}{t} + \nabla\cdot(uD) & = 0,
\end{align}
where $u$ is the velocity, $D$ is the layer depth, $b$ is the height
of the bottom of the fluid layer, $g$ is the acceleration due to
gravity, $u^{\perp} = \hat{k}\times u$, $\hat{k}$ is the outward
pointing normal to the surface (e.g. on the sphere), and $\nabla$ is
the restriction of the gradient tangential to the surface.

These equations have a Hamiltonian variational form
\begin{align}
  \langle w, u_t \rangle +
  \left\langle w/D, {\dede{H}{u}}^\perp (\nabla\cdot u^\perp + f)
  \right\rangle - \left\langle
  \nabla\cdot w, \dede{H}{D}
  \right\rangle & = 0, \quad \forall w,
  \\
  \left\langle \phi, D_t + \nabla\cdot \dede{H}{u} \right\rangle & = 0,
  \quad \forall \phi,
\end{align}
being deliberately vague about the spaces, with
\begin{align}
  H &= \int_\Omega \frac{D|u|^2}{2} + gD\left(\frac{D}{2} + b\right) \diff x,
\end{align}
hence
\begin{equation}
  \dede{H}{u} = Du, \quad \dede{H}{D} = \frac{|u|^2}{2} + g(D+b).
\end{equation}
Conservation of $H$ follows from
\begin{align}
  \label{eq:econ}
  \dot{H} = \langle \dede{H}{u}, u_t \rangle + \langle \dede{H}{D}, D_t \rangle,
\end{align}
and then choosing $w=\dede{H}{u}$, $\phi=\dede{H}{D}$ in the variational
formulation of the equations.

If we choose finite element spaces $V\in H_{\mathrm{div}}$, $Q\in L^2$,
then $\nabla\cdot u^\perp$ is not well-defined, hence we use a DG-type
formulation for that term, leading to the formulation seeking
$(u,D)\in V\times Q$ such that
\begin{align}
  \label{eq:u}
  \langle w, u_t \rangle -
  \left\langle \nabla^\perp_h\left(\frac{w}{D}\cdot{\dede{H}{u}}^\perp\right)
  u
  \right\rangle
  +
  \left\llangle 2\left\{n^\perp\left(\frac{w}{D}\cdot{\dede{H}{u}}^\perp\right)\right\},
  \tilde{u}
  \right\rrangle  
    +
  \left\langle w/D, {\dede{H}{u}}^\perp f
  \right\rangle - \left\langle
  \nabla\cdot w, \dede{H}{D}
  \right\rangle & = 0, \quad \forall w \in V,
  \\
  \left\langle \phi, D_t + \nabla\cdot \dede{H}{u} \right\rangle & = 0,
  \quad \forall \phi \in Q,
  \label{eq:D}
\end{align}
where $\nabla_h$ is the elementwise broken derivative,
$\llangle\cdot,\cdot\rrangle$ is the inner product integrating over
all the mesh facets, $\{\cdot\}$ indicates an average over the two
sides of a facet, and $\tilde{u}$ indicates the upwind value of $u$
(we can also use the average but the upwinding does not break energy
conservation and reduces gridscale oscillations). As before,
energy conservation follows from \eqref{eq:econ}, choosing
and then choosing $w=\dede{H}{u} \in V$, $\phi=\dede{H}{D}\in Q$
with variational derivatives now taken within $V\times Q$. [state
  and prove as a theorem]

To obtain an energy conserving time integration scheme we use a
continuous time Petrov-Galerkin scheme, with $(u,v) \in (V\otimes P,
Q\otimes P)$, $(w,\phi) \in (V\otimes P', Q\otimes P')$, where $P$ is
a degree $k$ polynomial space in time, and $P'$ is the space one
degree lower. A straightforward application of the scheme gives
\begin{align}
  \int_0^{\Delta t}\langle w, u_t \rangle -
  \left\langle \nabla^\perp_h\left(\frac{w}{D}\cdot{\dede{H}{u}}^\perp\right)
  u
  \right\rangle
  +
  \left\llangle 2\left\{n^\perp\left(\frac{w}{D}\cdot{\dede{H}{u}}^\perp\right)\right\},
  \tilde{u}
  \right\rrangle  
    +
  \left\langle w/D, {\dede{H}{u}}^\perp f
  \right\rangle - \left\langle
  \nabla\cdot w, \dede{H}{D}
  \right\rangle\diff t & = 0, \quad \forall w \in V\otimes P',
  \\
  \int_0^{\Delta t}\left\langle \phi, D_t + \nabla\cdot \dede{H}{u} \right\rangle \diff t & = 0,
  \quad \forall \phi \in Q\otimes P'.
\end{align}
The energy conservation does not work, because we can't choose
$w=\dede{H}{u}$, because $w \in V\times P'$ and $\dede{H}{u}\in V\times P$.
Hence, we make the modification
\begin{align}
  \label{eq:discrete u}
  \int_0^{\Delta t}\langle w, u_t \rangle -
  \left\langle \nabla^\perp_h\left(\frac{w}{D}\cdot{
    \mathbb{P}(\dede{H}{u})}^\perp\right)
  u
  \right\rangle
  +
  \left\llangle 2\left\{n^\perp\left(\frac{w}{D}\cdot{
    (\mathbb{P}\dede{H}{u}})^\perp\right)\right\},
  \tilde{u}
  \right\rrangle  
    +
  \left\langle w/D, {\mathbb{P}(\dede{H}{u})}^\perp f
  \right\rangle - \left\langle
  \nabla\cdot w, \dede{H}{D}
  \right\rangle\diff t & = 0, \quad \forall w \in V\otimes P',
  \\
  \int_0^{\Delta t}\left\langle \phi, D_t + \nabla\cdot \dede{H}{u} \right\rangle \diff t & = 0,
  \quad \forall \phi \in Q\otimes P',
  \label{eq:discrete D}
\end{align}
where $\mathbb{P}$ is the $L^2$ time projection operator into $P'$.

Then, energy conservation follows from
\begin{align}
H^{n+1} - H^n
& = \int_0^T \dot{H} \diff t, \\
& = \int_0^T \langle \dede{H}{u}, u_t \rangle + \langle \dede{H}{D},
D_t \rangle \diff t, \\
& = \int_0^T \langle \mathbb{P}(\dede{H}{u}), u_t \rangle + \langle \mathbb{P}(\dede{H}{D}),
D_t \rangle \diff t,
\end{align}
which becomes zero after using $w=\mathbb{P}(\dede{H}{u})$,
$\phi=\mathbb{P}(\dede{H}{D})$ in the motion equations
(\ref{eq:discrete u}-\ref{eq:discrete D}).
[state and prove as a theorem]

This can be made into a regular continuous Petrov Galerkin time scheme
without $\mathbb{P}$ operators by introducing an additional variable
$G\in V\times P$ with $G_t = \mathbb{P}(\dede{H}{u})$, written in
space-time variational form as
\begin{equation}
  \int_0^\Delta t
  \langle \gamma, G_t - \dede{H}{u} \rangle = 0, \quad
  \forall \gamma \in V\times P'.
\end{equation}
Then an equivalent formulation seeks
$(u, G, D) \in (V\otimes P)\times (V\otimes P)\times (Q\otimes P)$
such that
\begin{align}
  \label{eq:discrete u 2}
  \int_0^{\Delta t}\langle w, u_t \rangle -
  \left\langle \nabla^\perp_h\left(\frac{w}{D}\cdot
    G_t^\perp\right)
  u
  \right\rangle
  +
  \left\llangle 2\left\{n^\perp\left(\frac{w}{D}\cdot G_t^\perp\right)\right\},
  \tilde{u}
  \right\rrangle  
    +
  \left\langle w/D, G_t^\perp f
  \right\rangle - \left\langle
  \nabla\cdot w, \dede{H}{D}
  \right\rangle\diff t & = 0, \quad \forall w \in V\otimes P',
  \\
  \int_0^{\Delta t}\left\langle G_t - \dede{H}{u}, \gamma \right\rangle
  \diff t & = 0, \quad \forall \gamma \in V\otimes P', \\
  \int_0^{\Delta t}\left\langle \phi, D_t + \nabla\cdot \dede{H}{u} \right\rangle \diff t & = 0,
  \quad \forall \phi \in Q\otimes P'.
  \label{eq:discrete D 2}
\end{align}
This makes it possible to use Irksome's time Galerkin integrator without
modification.

Finally, does that the energy conservation proof still works
if we replace the time integral by a quadrature rule provided that
it is sufficiently high degree to compute
\[
\int_0^{\Delta t} \langle w, u_t \rangle \diff t
\]
without error? Similarly for the spatial inner products (thus enabling
nonexact integration due to division by $D$, upwinding of $u$, etc).
[State and prove as a theorem - I'm slightly worried about whether
  this still works in the time part with the $G_t$ modification. I
  think it doesn't, because reducing the degree to 1 in the quadrature
  causes energy errors that seem bigger than solver tolerance
  stuff). But we can discuss the quadrature for the space terms
  anyway. It also looks like we don't need to integrate the
  $\dede{H}{D}$ term exactly, which isn't a concern here but would be
  for ideal gas dynamics where the internal energy has fractional
  powers in it].

The novelty here is:
\begin{enumerate}
  \item we fully formulate and demonstrate energy conservation for
    high order time integrators,
  \item The equivalent formulation that works as a standard
    space-time scheme without projections,
  \item we have iterative schemes for solving (\ref{eq:discrete u 2}-
    \ref{eq:discrete D 2}) in a scalable manner. Previously people
    just did Picard iteration which doesn't really work.
\end{enumerate}

\subsection{SFLT}
Extension to SFLT, just add a noise term to $u$ in $\nabla \cdot u^\perp$
(and modify spatial discretisation of that term appropriately).
[We need to check with Darryl that this is the right thing.]

\section{Local energy conservation}

To develop a local energy conserving formulation, we introduce
the broken space $\tilde{V}\supset V$ consisting of the same local
shape functions but without continuity constraints, and the trace
space $T$, supported on the skeleton of the mesh, with
$u \in V\implies u\cdot n \in T$, and $\forall \gamma \in T$ $\exists
u \in V$ such that $u\cdot n = \gamma$.

We then have the broken/trace formulation finding $(u,D, \lambda)\in
\tilde{V}\times Q \times T$ such that
\begin{align}
  \nonumber
  \langle w, u_t \rangle -
  \left\langle \nabla^\perp_h\left(\frac{w}{D}\cdot{\dede{H}{u}}^\perp\right)
  u
  \right\rangle
    +
  \left\langle w/D, {\dede{H}{u}}^\perp f
  \right\rangle - \left\langle
  \nabla_h\cdot w, \dede{H}{D}
  \right\rangle & \\
  \label{eq:br u}
  \qquad +   
  \left\llangle 2\left\{n^\perp\left(\frac{w}{D}\cdot{\dede{H}{u}}^\perp\right)\right\},
  \tilde{u}
  \right\rrangle
  + \left\llangle \jump{w}, \lambda \right\rrangle
& = 0, \quad \forall w \in \tilde{V},
  \\
  \left\langle \phi, D_t + \nabla_h\cdot \dede{H}{u} \right\rangle & = 0,
  \quad \forall \phi \in Q, \\
  \left\llangle \jump{u}, \mu \right \rrangle & = 0,
  \quad \forall \mu \in T, \label{eq:br lambda}
\end{align}
where $\jump{u}=u^+\cdot n^+ + u^-\cdot n^-$.

The solveability of this system needs to be established, but it is
known for the linear rotating SWE. Assuming we have a solution, then
we automatically have a solution of the original unbroken system.
First, from the jump condition \eqref{eq:br lambda} we have that
$\tilde{u}\in V \subset \tilde{V}$. Second, taking $w \in \tilde{V}\subset
V$ in \eqref{eq:br u} means that $\jump{w}=0$, and we recover 
\eqref{eq:u}.

Further, any solution of (\ref{eq:u}-\ref{eq:D}) solves equations
(\ref{eq:br u}-\ref{eq:br lambda}). To see this, we define the operator
$L:T\to \tilde{V}$ by seeking $v, \psi \in \tilde{V}\times Q$ such that
\begin{align}
  \langle w, v \rangle - \langle \nabla_h\cdot w, \psi \rangle
  & = - \llangle \jump{w}, \lambda \rrangle, \quad \forall w \in \tilde{V}, \\
  \langle \phi, \nabla\cdot v \rangle & = 0, \quad \forall \phi \in Q,
\end{align}
which is solveable by standard mixed finite element arguments, and we
take $L\lambda=v$.

We claim that any $w \in \tilde{V}$ can be written as
$\bar{w}+L\lambda$ where $\bar{w}\in V$, for some $\lambda \in
T$. This is done by first defining $\lambda=\jump{w}$. Then,
$\bar{w}=w-L\lambda$ satisfies
\[
\jump{\bar{w}} = \jump{w} - \jump{L\lambda}
\]



Then, for any $\mu \in T$, we can use $w=L\mu$ in \eqref{eq:br u} to
give
\begin{align}
\nonumber
  \left\llangle \mu, \lambda \right\rrangle
  & = 
  -\langle L\mu, u_t \rangle +
  \left\langle \nabla^\perp_h\left(\frac{L\mu}{D}\cdot{\dede{H}{u}}^\perp\right)
  u
  \right\rangle
    -
  \left\langle L\mu/D, {\dede{H}{u}}^\perp f
  \right\rangle - \left\langle
  \nabla\cdot L\mu, \dede{H}{D}
  \right\rangle \\
  & \qquad -   
  \left\llangle 2\left\{n^\perp\left(\frac{L\mu}{D}\cdot{\dede{H}{u}}^\perp\right)\right\},
  \tilde{u}
  \right\rrangle, \quad \forall \mu \in T.  \label{eq:recon lambda}
  \end{align}

\end{document}

